\subsection{紧算子的谱}

引理 1. --- 设 E 是 R 上维数 \(\geqslant 2\) 的分离拓扑向量空间，X 是 E 的一个可数子集。补集 E - X 是连通的。

首先假设 E 的维数为 2。我们可以假设 \(E = R^{2}\) 赋有欧几里得范数（EVT, I, p. 14, 定理 2）。由于 X 可数，存在一个实数 \(r \in R_{+}^{*}\) ，使得以 0 为中心、r 为半径的圆 C 不与 X 相交。设 \(x \in E - X\) ；若 \(x \notin C\) ，存在一点 \(y \in C\) ，使得连接 x 与 y 的直线 \(L_{x}\) 不与 X 相交，因为 X 可数。集 E−X 是 C（它是连通的）与诸连通集 \(L_{x} \cup C\) （\(x \in \mathrm{E} - (\mathrm{X} \cup \mathrm{C})\) ）的并；这些集都包含 C，因此 E-X 是连通的（TG, I, p. 81, 命题 2）。

考察一般情形。对元素 \(x \in E - X\) 把 X 换成 X - x，我们化归到 \(0 \notin X\) 的情形。由于 E-X 是诸集 \(F - (X \cap F)\) （当 F 取遍 E 的 2 维子空间时所成的集）的并，且由前面的情形这些集都连通并包含 0，故集 E-X 是连通的（同前）。

引理 2. --- 设 S \(\subset \mathbf{C}\) 是 \(\mathbf{C} - \{0\}\) 中的一个无限的、离散的、有界的且闭的集。则 S 是某个序列 ( \(\lambda_{n}\) ) \(_{n\in N}\) 的值集，该序列由两两不同的非零复数组成，且序列 ( \(|\lambda_{n}|\) ) \(_{n\in N}\) 递减并收敛到 0。

对每个整数 \(i \geqslant 1\) ，集 \(\mathbf{A}_i\) （由复数 \(\lambda \in \mathrm{S}\) （满足 \(|\lambda| \geqslant \frac{1}{i}\) ）组成）在 \(\mathbf{C}\) 中是紧且离散的，故有限。用 \(a_i\) 表示它的基数。由于 S 无限，序列 \((a_i)\) 趋于 \(+\infty\) 。令 \(\mathbf{A}_0 = \varnothing\) ，\(a_0 = 0\) 。对每个 \(i \geqslant 1\) ，选取一个双射 \(n \mapsto \lambda_n\) （从区间 \([a_{i-1}, a_i[\) （\(\mathbf{N}\) 的）到 \(\mathbf{A}_i - \mathbf{A}_{i-1}\) 上），使得映射 \(n \mapsto |\lambda_n|\) 在 \([a_{i-1}, a_i[\) 上递减。序列 \((\lambda_n)_{n \in \mathbf{N}}\) 具有所要求的性质。

设 E 是无限维的完备可赋范空间。代数 \(\mathcal{L}^{\mathrm{c}}(\mathrm{E})\) 是 \(\mathcal{L}(\mathrm{E})\) 的一个不含单位元的子代数。回忆：对每个紧自同态 \(u \in \mathcal{L}^{\mathrm{c}}(\mathrm{E})\) ，谱 \(\mathrm{Sp}'_{\mathcal{L}^{\mathrm{c}}(\mathrm{E})}(u)\) 是 \(u\) 相对于含单位元子代数 \(\mathcal{L}^{\mathrm{c}}(\mathrm{E}) \oplus \mathbf{C1}_{\mathrm{E}}\) （\(\mathcal{L}(\mathrm{E})\) 的）的谱（I, p. 4, 第 4 目）。

命题 5. — 设 E 是完备可赋范空间，u 是 E 的一个紧自同态。\(\mathrm{Sp}_{\mathrm{s}}(u)\) 的每个元素都是有限谱重数的特征值。此外：

\begin{enumerate}
\def\labelenumi{\alph{enumi})}
\item
  若 E 有限维，则 \(\operatorname{Sp}(u) = \operatorname{Sp}_{\mathrm{s}}(u)\) ；
\item
  若 E 无限维，则 \(\mathrm{Sp}_{\mathrm{s}}(u) = \mathrm{Sp}(u) - \{0\}\) 且 \(\mathrm{Sp}_{\omega}(u) = \{0\}\) ；
\item
  若 \(\mathrm{Sp}_{\mathrm{s}}(u)\) 无限，则它是某个序列 \((\lambda_{n})_{n\in\mathbf{N}}\) 的值集，该序列由两两不同的非零复数组成，且序列 \((|\lambda_{n}|)_{n\in\mathbf{N}}\) 递减并收敛到 0；
\item
  若 E 无限维，则 \(\mathrm{Sp}(u) = \mathrm{Sp}'_{\mathcal{L}^c (\mathrm{E})}(u)\) 。
\end{enumerate}

\(\mathrm{Sp}_{\mathrm{s}}(u)\) 的每个元素都是有限谱重数的特征值（III, p. 83 的命题 1）。

断言 \(a\) 是初等的（III, p. 85, 例 1）。现在假设 E 无限维。

设 \(\lambda \in \mathrm{Sp}(u) - \{0\}\) 。则 \(u - \lambda 1_{\mathrm{E}}\) 是 E 的里斯自同态（III, p. 75 的命题 2 的推论 2），故 \(\lambda\) 属于 \(\mathrm{Sp}_{\mathrm{s}}(u)\) 。若 E 无限维，则 \(\mathrm{Sp}_{\omega}(u)\) 非空（III, p. 87, 定理 1, a)）。必然有 \(\mathrm{Sp}_{\omega}(u) = \{0\}\) ，由此得 \(b\) 。

集 \(\mathrm{Sp}_{\mathrm{s}}(u)\) 是离散且有界的。此外，我们有 \(\mathrm{Sp}_{\mathrm{s}}(u) = \mathrm{Sp}(u) \cap (\mathbf{C} - \{0\})\) （依照 \(b\) ），故 \(\mathrm{Sp}_{\mathrm{s}}(u)\) 在 \(\mathbf{C} - \{0\}\) 中是闭的。于是断言 \(c\) 由引理 2 得出。

\(u\) 的谱是可数的，故其在 \(\mathbf{C}\) 中的补集是连通的（引理 1）。把 I, p. 28 的命题 6 的推论应用于含单位元子代数 \(\mathcal{L}^{\mathrm{c}}(\mathrm{E})\oplus \mathbf{C}1_{\mathrm{E}}\) （\(\mathcal{L}(\mathrm{E})\) 的），我们断定 \(\mathrm{Sp}(u) = \mathrm{Sp}'_{\mathcal{L}^{\mathrm{c}}(\mathrm{E})}(u)\) 。

命题 6. --- 设 E 是 C 上的完备可赋范空间，u 是 E 的一个紧自同态。

\begin{enumerate}
\def\labelenumi{\alph{enumi})}
\item
  设 \(f \in \mathcal{O}(\mathrm{Sp}(u))\) 满足 \(f(0) = 0\) 。自同态 \(f(u)\) 是紧的；
\item
  假设 E 是复希尔伯特空间且 u 是正规的。设 f 是 \(\mathrm{Sp}(u)\) 上的一个连续函数，满足 \(f(0)=0\) 。正规自同态 \(f(u)\) 是紧的。
\end{enumerate}

此外，若 E 无限维，则逆命题成立；若 E 有限维，则条件 \(f(0) = 0\) 不是必要的。

我们可以假设 E 无限维。

我们来证明 \(a\) 及其逆。自同态 \(u\) 是巴拿赫代数 \(\mathcal{L}^{\mathrm{c}}(\mathrm{E})\) 的一个元素。由于 E 无限维，我们有等式 \(\mathrm{Sp}(u) = \mathrm{Sp}'_{\mathcal{L}^{\mathrm{c}}(\mathrm{E})}(u)\) （由命题 5 的 \(d\) ）。全纯泛函演算的元素 \(f(u)\) （该演算是把 \(\mathcal{L}^{\mathrm{c}}(\mathrm{E})\) 添加单位元后得到的含单位元巴拿赫代数的全纯泛函演算）属于 \(\mathcal{L}^{\mathrm{c}}(\mathrm{E})\) 当且仅当 \(f(0) = 0\) （I, p. 88）。而且这个元素与元素 \(f(u)\) （\(\mathcal{L}(\mathrm{E})\) 的）重合（I, p. 75 的命题 7），故 \(f(u)\) 是紧的当且仅当 \(f(0) = 0\) 。

断言 \(b\) 及其逆的证明完全相同，只需考察 C*-代数 \(\mathcal{L}^{\mathrm{c}}(\mathrm{E})\) 的连续泛函演算（I, p. 110, 定义 5）。

推论. --- 设 E 和 F 是希尔伯特空间，u 是从 E 到 F 中的一个连续线性映射。线性映射 u 是紧的，当且仅当 E 的自同态 \(|u|\) 是紧的。

设 \((j, |u|)\) 是 u 的极分解（I, p. 140 的定义 4）。由于 \(u = j|u|\) 且 \(|u| = \sqrt{u^{*}u}\) （I, p. 139 的命题 10），该等价性由 III, p. 5 的命题 3 以及前一命题的断言 b) 得出。

